%%%%%%%%%%%%%%%%%%%%%%%%%%%%%%%%%%%%%%%%%%%%%%%%%%%%%%%%%%%%%
\section{Numerical methods}\label{sec:numerical_methods}
%%%%%%%%%%%%%%%%%%%%%%%%%%%%%%%%%%%%%%%%%%%%%%%%%%%%%%%%%%%%%
% Every result in Sections 4-6 was produced with the partitioned procedures
% described here (solids4foam fluidSolidInterface: fixedRelaxation, Aitken,
% IQNILS; Robin condition elasticWallPressure). The experimental
% quasi-monolithic solver of the parent paper is not used; its description is
% preserved, outside the build, in sectionNumericalModelQuasiMonolithic.tex.
% Source checked: solids4foam development (src/solids4FoamModels:
% fluidModels/pimpleFluid, solidModels, fluidSolidInterfaces) and the case
% dictionaries of the tutorials in Section 5.

All results in this paper were obtained with the partitioned fluid--solid interaction framework of the open-source toolbox solids4foam~\citep{Cardiff2025}, which builds on the OpenFOAM finite volume library~\citep{Weller1998} and follows the partitioned approach of \citet{Tukovic2018}.
The fluid and the solid are discretised on separate cell-centred finite volume meshes and solved by separate field solvers, which exchange interface displacements and tractions within each time step.
This section summarises the components that determine the discretisation, temporal and iterative errors studied in Sections~\ref{sec:verification_methodology}--\ref{sec:cross_benchmark}; implementation detail is in the solids4foam source and in the case dictionaries that accompany each benchmark.
Case-specific choices (time scheme, convection scheme, solid variant, coupling algorithm and tolerances) are stated with each benchmark in Section~\ref{sec:benchmark_suite}.

%%--------------------------------------------------------------------------------------------------------------------%%
\subsection{Fluid}\label{sec:fluid_solver}
%%--------------------------------------------------------------------------------------------------------------------%%

The incompressible Navier--Stokes equations in ALE form (Section~\ref{sec:fluid_governing_eqn}) are solved by the \texttt{pimpleFluid} model, a solids4foam port of the OpenFOAM \texttt{pimpleFoam} solver.
Velocity and pressure are collocated at cell centres; the face fluxes are obtained by Rhie--Chow-type momentum interpolation~\citep{Rhie1983}, and the pressure--velocity coupling is the segregated PIMPLE algorithm (PISO pressure correctors within SIMPLE-type outer correctors), with the numbers of outer and pressure correctors stated per case.
Spatial operators are the standard second-order finite volume operators: midpoint quadrature, linear face interpolation, least-squares or Gauss cell gradients, and Laplacian face fluxes with explicit non-orthogonal correction~\citep{Jasak1996}.
Convection is discretised with central (\texttt{linear}) or limited linear-upwind interpolation, depending on the case.
The time derivative is discretised with the second-order backward difference formula (BDF2, \texttt{backward}) unless stated; implicit Euler is used only in the deliberate time-scheme comparison of Section~\ref{sec:3dtube}.
On the moving mesh the convective flux is the relative flux $\bb{\Gamma}_f\cdot(\bb{v}_f - \bb{v}^\Gamma_f)$, where the mesh flux $\bb{\Gamma}_f\cdot\bb{v}^\Gamma_f$ is the volume swept by face $f$.
With BDF2 the swept volumes of the last two steps are combined with the BDF2 weights, so that the discrete space-conservation law (Equation~\eqref{eqn:gcl}) holds for the backward volume derivative.
On the interface the fluid velocity equals the velocity of the moving wall, and the interface traction passed to the solid is the sum of the pressure and viscous tractions on the fluid interface faces.

\subsection{Fluid mesh motion}\label{sec:fluid_mesh_motion}

The interior fluid mesh follows the interface through a Laplacian equation for the cell motion velocity (\texttt{velocityLaplacian}), whose diffusivity varies with the inverse (in most cases the inverse square) of the distance to the moving boundaries, so that cells near the interface move almost rigidly~\citep{Tukovic2018}.
The interface displacement is the solid displacement transferred to the fluid interface (Section~\ref{sec:interface_transfer}); the point motion is obtained from the cell motion velocity by interpolation, and the mesh points are moved once per coupling iteration.
Mesh motion is a numerical device: it changes the discrete solution only through the mesh quality and through the space-conservation and moving-wall treatments described above.

%%--------------------------------------------------------------------------------------------------------------------%%
\subsection{Solid}\label{sec:solid_solver}
%%--------------------------------------------------------------------------------------------------------------------%%

The solid momentum equation (Section~\ref{sec:solid_model}) is discretised by a cell-centred finite volume method with the total displacement as unknown~\citep{Jasak2000, Cardiff2021}.
Two kinematic formulations are used: a small-strain (linear geometry) formulation for the linear-elastic walls (ringAddedMass, womersleyTube, 3dTube), and a total Lagrangian formulation for the finite-strain problems, with St Venant--Kirchhoff (Hron--Turek, collapsibleChannel, mokLidDrivenCavity) or neo-Hookean (Hessenthaler) material laws.
The time derivative uses BDF2, and each solid solve is a Newton--Krylov iteration through the PETSc SNES interface~\citep{PETSc, Cardiff2026}.
Two spatial variants are used.
The \emph{standard} variant uses a linear least-squares reconstruction of the displacement and a compact-stencil momentum stabilisation, a Rhie--Chow-type difference between the compact and the reconstructed face normal gradient, scaled by a user factor~\citep{Cardiff2026}; it is formally second order, but the stabilisation is not negligible on coarse meshes and stiffens thin, bending-dominated structures (Section~\ref{sec:fsi1}).
The \emph{high-order} variant evaluates the face fluxes of the residual from a cubic moving-least-squares reconstruction of the displacement~\citep{Batistic2026}.
Which variant is used is stated per case, because for thin bending-dominated walls the choice can dominate the coupled error (Section~\ref{sec:collapsible}).

%%--------------------------------------------------------------------------------------------------------------------%%
\subsection{Interface transfer}\label{sec:interface_transfer}
%%--------------------------------------------------------------------------------------------------------------------%%

Displacements are transferred from the solid to the fluid interface at the interface points, and tractions from the fluid to the solid interface at the faces.
Where the two interface meshes coincide (ringAddedMass, womersleyTube, collapsibleChannel, mokLidDrivenCavity), a direct one-to-one map is used.
Where they do not (Hron--Turek, 3dTube, Hessenthaler), face tractions are transferred conservatively with arbitrary-mesh-interface (AMI) weights, and point displacements by barycentric interpolation on a triangulation of the donor interface faces.
The interface transfer is part of the discretisation: it is an operator that must converge with refinement like any other, and in 3dTube a defect in it changed with mesh level (Section~\ref{sec:3dtube}).

%%--------------------------------------------------------------------------------------------------------------------%%
\subsection{Coupling algorithms}\label{sec:coupling_algorithms}
%%--------------------------------------------------------------------------------------------------------------------%%

Strong coupling is enforced by outer (coupling) iterations within each time step.
Two interface formulations are used.

\paragraph{Dirichlet--Neumann.}
The fluid receives the interface displacement (mesh motion and wall velocity) and the solid receives the fluid traction.
Each time step starts with a solid predictor, i.e.\ a solid solve with the fluid traction of the previous step, whose displacement is passed to the fluid without relaxation in the first iteration~\citep{Tukovic2018}.
Subsequent interface displacements are updated by fixed under-relaxation, by Aitken dynamic relaxation~\citep{kuttler2008fixed}, or by the interface quasi-Newton inverse least-squares method (IQN-ILS)~\citep{Degroote2009}, which builds an approximate inverse Jacobian of the interface residual from the secant information of the current and, optionally, previous time steps; secant columns that are repeated or nearly linearly dependent are filtered before the least-squares solve.

\paragraph{Robin--Neumann.}
The fluid pressure satisfies a Robin condition on the interface that accounts for the solid inertia through an estimate $\rho_S h_S$ of the solid interface impedance~\citep{Tukovic2018b}, and the iterations are unrelaxed fixed-point iterations.
The Robin condition modifies the discrete interface condition within the iterations; at convergence the interface flux is made consistent with the mesh motion.
Because it changes the iteration rather than only accelerating it, the converged Robin--Neumann and Dirichlet--Neumann solutions need not agree to the coupling tolerance (Section~\ref{sec:ref_hierarchy}).

\paragraph{Convergence test.}
The interface residual $\bb{r}$ is the difference between the solid interface point positions, transferred to the fluid interface, and the fluid interface point positions.
Two normalised measures are formed,
\begin{equation}
	r_1 = \frac{\|\bb{r}\|}{\max_k \|\bb{r}^{(k)}\|},
	\qquad
	r_2 = \frac{\|\bb{r}\|}{\|\bb{d}^\Gamma\|_{\max}},
	\label{eq:coupling_residuals}
\end{equation}
where the maximum in $r_1$ is over the coupling iterations $k$ of the current step and $\|\bb{d}^\Gamma\|_{\max}$ is the largest norm of the \emph{total} interface displacement seen so far.
In the legacy default a step is accepted when $\min(r_1, r_2)$ is below the coupling tolerance; a strict option requires $\max(r_1, r_2)$ to be below it.
The two tests differ greatly in slowly varying phases with a large total interface displacement, where $r_2$ can be small after a single relaxed iteration whatever the reduction of the residual within the step (Section~\ref{sec:iterative_errors}).
Robin--Neumann iterations additionally require the change of interface pressure and the interface leakage flux to fall below their own tolerances.
By default, a step that does not meet the test within the maximum number of iterations stops the run; the default test can nevertheless accept insufficiently converged steps, as shown in Section~\ref{sec:iterative_errors}.

%%--------------------------------------------------------------------------------------------------------------------%%
\subsection{Solution procedure and toolchain}\label{sec:implementation_details}
%%--------------------------------------------------------------------------------------------------------------------%%

The partitioned procedure for one time step is summarised in Algorithm~\ref{alg:partitioned_fsi}.

\begin{algorithm}[htbp]
\caption{Partitioned FSI procedure for one time step (solids4foam \texttt{fluidSolidInterface}).}
\label{alg:partitioned_fsi}
\begin{algorithmic}[1]
\State Advance time; solve the solid with the fluid traction of the previous step (predictor).
\State Transfer the predicted interface displacement to the fluid interface.
\State For coupling iterations $k = 1, 2, \ldots$, up to the maximum number:
\State \hspace{1.5em}move the fluid mesh to the current interface displacement and update the mesh fluxes;
\State \hspace{1.5em}solve the fluid (PIMPLE outer and pressure correctors) with the moving-wall velocity, or with the Robin pressure condition;
\State \hspace{1.5em}transfer the fluid interface traction to the solid and solve the solid (Newton--Krylov);
\State \hspace{1.5em}transfer the solid interface displacement to the fluid interface and evaluate $\bb{r}$, $r_1$ and $r_2$;
\State \hspace{1.5em}stop if the convergence test (and, for Robin--Neumann, the pressure and flux tests) is met;
\State \hspace{1.5em}otherwise update the interface displacement: fixed relaxation, Aitken, IQN-ILS, or none (Robin--Neumann).
\end{algorithmic}
\end{algorithm}

The calculations used the OpenCFD releases OpenFOAM-v2412 and OpenFOAM-v2512 with PETSc, on Linux (x86\_64) and, for some of the inexpensive cases, macOS; the release, platform and solids4foam commit are stated with each result.
The compiler and its optimisation level are recorded as part of the configuration: the default OpenFOAM build rules compile at \texttt{-O3}, whereas the EasyBuild toolchain used on one of the clusters compiles at \texttt{-O2}, and this difference exposed a solver defect (Section~\ref{sec:3dtube}).
\todo{State per case the exact solids4foam commit, OpenFOAM release, compiler and optimisation flags, PETSc version and platform (see Section~\ref{sec:qoi} and paper-data/manuscript\_update\_2026-10.md). Confirm that the predictor/relaxation description matches the commit used for each result (the \texttt{predictor} default changed in solids4foam PR \#492).}
